% part: intuitionistic-logic
% chapter: semantics
% section: topological-semantics

\documentclass[../../../include/open-logic-section]{subfiles}

\begin{document}

\olfileid[id]{int}{sem}{top}

\olsection{Semantik Topologis}

Cara lain untuk memberikan semantik bagi logika intuisionistik ialah dengan
menggunakan konsep matematis topologi.

\begin{defn}
  Misalkan $X$ merupakan himpunan. Suatu \emph{topologi pada~$X$} ialah
  himpunan $\Top{O}\subseteq\Pow{X}$ yang memenuhi sifat-sifat di bawah ini.
  !!^{element}s dari~$\Top{O}$ disebut \emph{himpunan terbuka} dalam topologi
  tersebut. Himpunan $X$ bersama $\Top{O}$ disebut \emph{ruang topologis}.
  \begin{enumerate}
  \item Himpunan kosong dan seluruh ruang bersifat terbuka:
    $\emptyset$, $X\in\Top{O}$.
  \item Himpunan terbuka tertutup terhadap irisan berhingga: jika
    $U$, $V\in\Top{O}$, maka $U\cap V\in\Top{O}$
  \item Himpunan terbuka tertutup terhadap gabungan sebarang: jika
    $U_i\in\Top{O}$ untuk semua $i\in I$, maka
    $\bigcup\Setabs{U_i}{i\in I}\in\Top{O}$.
  \end{enumerate}
\end{defn}

Kita dapat menulis $X$ bagi suatu ruang topologis jika koleksi himpunan
terbukanya dapat disimpulkan dari konteks. Namun, perlu diperhatikan bahwa
$X$ baru dapat disebut ruang topologis setelah dilengkapi dengan koleksi
himpunan terbuka.

\begin{defn}
  Suatu \emph{model topologis} logika proposisional intuisionistik ialah
  tripel $\mModel{X}=\tuple{X,\Top{O},V}$, dengan $\Top{O}$ merupakan topologi
  pada~$X$ dan $V$ merupakan fungsi yang menetapkan suatu himpunan terbuka
  dalam $\Top{O}$ kepada setiap variabel proposisional.

  Dengan diberikan model topologis~$\mModel{X}$, kita dapat mendefinisikan
  $\Prop{X}{!A}$ secara induktif sebagai berikut:
  \begin{enumerate}
  \item $\Prop{X}{\lfalse}=\emptyset$
  \item $\Prop{X}{p}=V(p)$
  \item $\Prop{X}{!A\land !B}=\Prop{X}{!A}\cap\Prop{X}{!B}$
  \item $\Prop{X}{!A\lor !B}=\Prop{X}{!A}\cup\Prop{X}{!B}$
  \item $\Prop{X}{!A\lif !B}=\Interior{(X\setminus\Prop{X}{!A})\cup
    \Prop{X}{!B}}$
  \end{enumerate}
  Di sini, $\Interior{V}$ merupakan fungsi yang memetakan himpunan
  $V\subseteq X$ ke \emph{interiornya}, yakni gabungan semua himpunan terbuka
  yang termuat di dalamnya. Dengan kata lain,
  \[
  \Interior{V}=\bigcup\Setabs{U}{U\subseteq V\text{ dan }U\in\Top{O}}.
  \]
\end{defn}

Perhatikan bahwa interior setiap himpunan selalu terbuka karena merupakan
gabungan himpunan-himpunan terbuka. Jadi, $\Prop{X}{!A}$ selalu merupakan
himpunan terbuka.

Walaupun semantik topologis sangat abstrak, ada cara memahaminya yang dapat
memberikan motivasi. Andaikan bahwa !!{element}s, atau ``titik-titik,'' dari
$X$ merupakan titik tempat pernyataan dapat dievaluasi. Himpunan semua titik
tempat $!A$ benar ialah proposisi yang dinyatakan oleh~$!A$. Tidak setiap
himpunan titik merupakan proposisi yang mungkin; hanya !!{element}s
dari~$\Top{O}$ yang demikian. $!A\Entails !B$ jika dan hanya jika $!B$ benar
pada setiap titik tempat~$!A$ benar, yakni
$\Prop{X}{!A}\subseteq\Prop{X}{!B}$, untuk semua~$X$. Pernyataan absurd
$\lfalse$ tidak pernah benar, sehingga $\Prop{X}{\lfalse}=\emptyset$.

Bagaimana proposisi yang dinyatakan oleh $!B\land !C$, $!B\lor !C$, dan
$!B\lif !C$ harus berhubungan dengan proposisi yang dinyatakan oleh $!B$
dan~$!C$ agar hukum-hukum yang valid secara intuisionistik berlaku, yakni agar
$!A\Proves !B$ jika dan hanya jika
$\Prop{X}{!A}\subseteq\Prop{X}{!B}$? Kita mensyaratkan
$\lfalse\Proves !A$ untuk setiap $!A$, dan hal ini terpenuhi karena
$\emptyset\subseteq U$ untuk semua $U$. Karena $!B\land !C\Proves !B$, kita
mensyaratkan $\Prop{X}{!B\land !C}\subseteq\Prop{X}{!B}$, dan serupa dengan
itu $\Prop{X}{!B\land !C}\subseteq\Prop{X}{!C}$. Himpunan terbesar yang
memenuhi $W\subseteq U$ dan $W\subseteq V$ ialah $U\cap V$. Sebaliknya,
$!B\Proves !B\lor !C$ dan $!C\Proves !B\lor !C$, sehingga kita mensyaratkan
$\Prop{X}{!B}\subseteq\Prop{X}{!B\lor !C}$ dan
$\Prop{X}{!C}\subseteq\Prop{X}{!B\lor !C}$. Himpunan terkecil~$W$ yang
memenuhi $U\subseteq W$ dan $V\subseteq W$ ialah $U\cup V$.

Definisi bagi $\lif$ lebih rumit: $!A\lif !B$ menyatakan proposisi terlemah
yang, jika digabungkan dengan $!A$, mengikutkan $!B$. Fakta bahwa
$!A\lif !B$ bersama $!A$ mengikutkan~$!B$ jelas dari
$(!A\lif !B)\land !A\Proves !B$. Jadi, $\Prop{X}{!A\lif !B}$ harus merupakan
himpunan terbuka terbesar sedemikian sehingga
$\Prop{X}{!A\lif !B}\cap\Prop{X}{!A}\subseteq\Prop{X}{!B}$; dari sinilah
definisi kita diperoleh.

\end{document}
